\section{System Architecture and Design}\label{sec:architecture}
Achieving real-time video segmentation requires a design that decouples computational execution from user interaction. This section details the modular pipeline that separates the classical and quantum-inspired engines and integrates them with an interactive UI.

\subsection{Asynchronous Threading Model and System Pipeline}
To maintain a responsive UI thread during heavy GPU calculations, the framework uses a dual-threaded asynchronous architecture \cite{williams2019c++}, as illustrated in Fig.~\ref{fig:system_architecture}.

\begin{figure}[!b]
\centering
\resizebox{\columnwidth}{!}{%
\begin{tikzpicture}[
    node distance=1.2cm and 1.5cm,
    block/.style={rectangle, draw, rounded corners=3, fill=blue!10, draw=blue!80, line width=1pt, minimum width=2.6cm, minimum height=1cm, align=center, font=\small},
    mutex/.style={rectangle, draw, rounded corners=3, fill=orange!10, draw=orange!80, line width=1pt, minimum width=2.6cm, minimum height=0.8cm, align=center, font=\small},
    gpu/.style={rectangle, draw, rounded corners=3, fill=green!10, draw=green!80, line width=1pt, minimum width=2.6cm, minimum height=1cm, align=center, font=\small},
    line/.style={-{Stealth[scale=1.2]}, line width=1.2pt},
    optline/.style={-{Stealth[scale=1.2]}, line width=1.2pt, dashed}
]
    % Layout
    \node[block] (UI) {UI Thread\\(Dear ImGui)};
    \node[mutex, below=of UI] (Mutex) {Mutex Lock\\(Shared State)};
    \node[block, below=of Mutex] (Worker) {Worker Thread\\(OpenCV Loop)};
    \node[gpu, right=of Worker] (GPU) {GPU Kernels\\(CUDA Engine)};

    % Arrows
    \draw[optline] (UI.south west) -- (Mutex.north west) node[midway, left, font=\footnotesize] {Update Config (Opt)};
    \draw[line] (Mutex.north east) -- (UI.south east) node[midway, right, font=\footnotesize] {Pull Textures};
    \draw[line] (Mutex) -- (Worker) node[midway, right, font=\footnotesize] {Read Config};
    \draw[line] (Worker) -- (GPU) node[midway, above, font=\footnotesize] {Launch};
    \draw[line] (GPU.north) |- (Mutex.east) node[pos=0.7, above, font=\footnotesize] {Commit Data};
\end{tikzpicture}%
}
\caption{Dual-Threaded Architecture with Asynchronous CUDA Pipeline Staging}
\label{fig:system_architecture}
\end{figure}

The UI Thread handles the window event loop and renders the immediate-mode Dear ImGui panels \cite{ImGui}. The Worker Thread runs the active clustering solver, looping through BGR frame acquisition via OpenCV \cite{Bradski2000}, CUDA-based distance calculation \cite{CudaGuide}, and direct GPU pixel re-coloring. This asynchronous boundary isolates computational latency from frame display pacing.

\subsection{Object-Oriented Design Patterns}
To maintain modularity, the system integrates several standard software engineering design patterns:
\begin{itemize}
    \item \textbf{Strategy Pattern:} Solvers are structured around the abstract \texttt{KMeansEngine} class. The \texttt{ClassicalEngine} and \texttt{QuantumEngine} act as interchangeable strategies, allowing live hot-swapping \cite{gamma1995design}.
    \item \textbf{Factory Method:} The \texttt{ClusteringFactory} instantiates the solver and centroid initialization variants based on configuration enums \cite{gamma1995design}.
    \item \textbf{Facade Pattern:} The \texttt{ClusteringManager} exposes a clean interface to hide raw preprocessing, kernel execution, and VRAM mapping operations \cite{gamma1995design}.
    \item \textbf{RAII (Resource Acquisition Is Initialization):} The \texttt{CudaAssignmentContext} controls memory allocation and release, ensuring device buffers are automatically freed during runtime layout adjustments \cite{stroustrup2022tour}.
    \item \textbf{Bridge Pattern (Pimpl):} Used in the manager classes as a compiler firewall to isolate CUDA and OpenCV includes from the global project compilation scope \cite{meyers2014effective}.
\end{itemize}

\subsection{Data Normalization and Preprocessing}
Pixels are converted into BGR features normalized to $[0, 1]$ and spatial features normalized by the frame dimensions: $\mathbf{x}_n = [B_n, G_n, R_n, \lambda_{\text{spatial}} x_n, \lambda_{\text{spatial}} y_n]^T$, with $\lambda_{\text{spatial}} = 10^{-6}$ \cite{shi2000normalized, achanta2012slic}. Centroid estimation is performed on a subsampled sparse manifold by stride sampling the frame at spatial stride $S \in [1, 16]$. By the Law of Large Numbers, this uniform sampling retains representation while reducing complexity by a factor of $S^2$ \cite{mouton2020stride}.

Thread communication is orchestrated through a thread-safe mutex bridge using \texttt{std::mutex} and \texttt{std::scoped\_lock} \cite{williams2019c++}, preventing race conditions when hot-swapping hyperparameters or transferring final textures (see Fig.~\ref{fig:system_architecture}).
